%% ma: L11-B16-0063
%% lop: 11
%% chuong: 5
%% bai: 16
%% dang: 11.16.4
%% loai: TLN
%% muc: VD
%% dap_an: "2"
%% nguon: "tai-lieu-hoc-tap-mon-toan-lop-11-theo-chuong-trinh-giao-duc-pho-thong-2018.pdf (Lớp Toán Thầy Duy, Phạm Lê Duy, THPT Chu Văn An), Chương 3 Bài 2, Phần 3, Câu 1 (tr.378)"
%% kien_thuc: "Lớp 11 Bài 16 (giới hạn của hàm số)"
%% trang_thai: chinh-thuc
%% ghi_chu: "Đề gốc dùng vế phải để chỉ dạng biến đổi; kho giải theo giá trị giới hạn (-1/2) và đáp số gốc 2 khớp. Đã duyệt 03/11/2026."
\begin{ex}
Giới hạn \[\lim\limits_{x\to+\infty}\left(\sqrt{x^2-x}-\sqrt[3]{x^3+1}\right)=\lim\limits_{x\to+\infty}\dfrac{-ax}{bx\left(\sqrt{1-\dfrac1x}+1\right)}\] với $a,b$ là các số tự nhiên và $\dfrac ab$ là phân số tối giản. Tính $P=a+b^2$.
\shortans{2}
\loigiai{\sloppy Giả sử vế phải có giới hạn $-\dfrac a{2b}$. Mặt khác $\sqrt{x^2-x}-x=\dfrac{-x}{\sqrt{x^2-x}+x}=\dfrac{-1}{\sqrt{1-\frac1x}+1}\to-\dfrac12$ và $\sqrt[3]{x^3+1}-x\to0$, nên giới hạn đã cho bằng $-\dfrac12$. Do đó $\dfrac ab=1$, tối giản nên $a=b=1$ và $P=1+1=2$.}
\end{ex}
